% Français 9115, batch 32 — Gallica views 621–640.
% Pass: Opus 5.5 (claude-opus-5-5), 2026-09-22 — first pass, unchecked against the views by a human.
%
% What the twenty views hold. Ten written pages on nine physical leaves,
% in her hand as far as the leaves show, in two pieces; none of it number
% theory. Batches 29–31 (views 561–620) were not on disk when this pass
% read its views; batch 31 appeared before this file was finished, and its
% header and last page were read then (see Seams). Batch 33, likewise
% written in parallel, continues with heat in the sphere and the cylinder.
%
% (1) v. 622–626: three pages of a French draft, heavily corrected, on the
% vibrations of cylindrical surfaces — square glass plates bent to
% different curvatures, compared with flat plates. Paginated at the head,
% by her, 2 (v. 622), 3 (v. 624), 4 (v. 626); page 1 is v. 620, outside
% this batch (glanced at only: an equation of the cylindrical surface of
% circular base, kb²{d⁴z/dx⁴ + 2d⁴z/dx²dy² + d⁴z/dy⁴} + (1/R²){…} = z″,
% and « le tems me manque … pour faire exécuter la totalité des
% experiences »). The three pages run on without a break:
%   v. 622  (2) she will give her experiments and compare the few facts
%           that lend themselves to it with the calculation in the
%           simplest case; N.o: to make the experiments comparable with
%           flat surfaces she had square plates cut and then bent; a
%           cancelled passage on cylindrical portions of arc A; before
%           bending she tried the tones of the simplest nodal figures,
%           and saw « avec étonnement » that Chladni's first figure, two
%           straight lines crossing at the centre, forms the same way on
%           the cylindrical surface as on the plate.
%   v. 624  (3) on some forty surfaces tried before and after bending, seven
%           or eight times a half-tone of difference, the bent surface
%           higher; she suspects the high temperature needed to make the
%           glass flexible, and « le peu d'activité des ateliers » at the
%           time kept her from varying the trials; N.o: on at least forty
%           square plates the interval between the tones of the figures
%           of two perpendicular lines and of three lines is the most
%           constant, differing only twice and by less than a half-tone; a
%           marginal note: on cylindrical surfaces this interval varies
%           with the curvature.
%   v. 626  (4) on square plates a given number of nodal lines always goes
%           with the same tone, the principle that guided Chladni's
%           classification of the distortions « et la théorie a
%           pleinement confirmé, comme je l'ai dit ailleurs »; on
%           cylindrical surfaces figures with the same number of lines
%           parallel to the curved or to the straight sides give different
%           tones; the figure of two diagonals, « suivant la théorie que
%           j'ai developpée N.o 11 de mon memoire pour 1813 », is
%           equivalent to that of two lines parallel to the sides, which
%           Chladni says never forms on square plates. The page ends at the
%           foot of the leaf on « a été expliquée ailleurs »; nothing after
%           it in this batch continues it.
% (2) v. 630–640: calculation leaves, clean or lightly corrected, on the
% propagation of heat in a sphere, a cone, a cylinder, a ring and a
% circular plate, in polar coordinates r, θ, φ, the temperature v, and a
% factor K/C.D (D possibly Ω, as batch 28 found at v. 549). Her own
% numbers here are 7 (v. 632), 8 (v. 635), 7 (v. 636), 6 (v. 639), and
% none on v. 630 and 640; the leaves are not bound in the order of those
% numbers, and the batch does not settle which order was hers.
%   v. 630  headed « cône »: the transformation of d²v/dx² + d²v/dy² +
%           d²v/dz² for the cone, taking derivatives with respect to
%           sin θ rather than θ, because sin θ measures the base of the
%           cone where θ measured the spherical cap; first derivatives,
%           then d²v/dx² expanded. The page stops at its foot.
%   v. 632  (7, struck) the planes through the vertex of the opposite cone
%           perpendicular to a given side, their intersections with the
%           planes zx and xy and their angles with the axis of x, written
%           as fractions Av/(π/2); a marginal « N.ota on a A + A′ = π/2 »;
%           small pen drawings of flowers and leaves down the left margin.
%   v. 635  (8) the cylinder, the ring (x constant in the equation of the
%           cylinder), then the cone, whose element ds is inclined to its
%           base, and the circular plate as the limit of an infinitely
%           flattened cone; a new paragraph begun and broken off on
%           « l'an ».
%   v. 639  (6, struck) the molecule of the sphere, sides dr, r sin θ dφ,
%           r dθ; the heat through its three faces and the opposite faces;
%           the equation dv/dt = K/C.D {d²v/dr² + (2/r)dv/dr + …}; then
%           « Si au lieu de considérer le solide conique qui s'appuie sur
%           la calotte sphérique on veut s'occuper du simple cône » — the
%           molecule of the cone. The sentence runs on to v. 636.
%   v. 636  (7, struck) the same for the molecule of the cone, sides dr,
%           r sin θ dφ, r cos θ dθ; the equation obtained differs from
%           « celle que j'ai trouvé d'une autre manière » by a factor
%           cos θ on the terms in d²v/(d sin θ)² and dv/d sin θ, explained
%           by the obliquity of the flux; a long marginal note continues
%           the text and ends « Retournez ».
%   v. 637  the back of that leaf: the molecule crossed obliquely; « alors
%           les équations seront identiques. »
%   v. 640  (dz/dx)² and d²z/dy² in the same coordinates, then
%           « Recherches sur les cones »: the element of the surface
%           parallel to the base, r′ = r cos θ / cos(θ − dθ), and the areas
%           crossed by the heat in the directions dv/dr and dv/dφ. The page
%           stops three quarters down; v. 641 is its blank back.
%
% Views not transcribed. v. 621, 623, 625, 627: the blank backs of the
% leaves of v. 620, 622, 624, 626, the writing showing through reversed
% (checked by mirroring v. 623 and 627). v. 628–629: the leaf inked 308,
% blank on both sides. v. 631, 633: blank backs of the leaves of v. 630
% and 632. v. 634: the side of the leaf inked 311 that carries only that
% number; its other side, v. 635, is written. v. 638: likewise the side of
% the leaf inked 313 with the number only; its other side is v. 639. No
% microfilm target and no repeated exposure in the batch.
%
% Seams. 620/621: view 620 (page 1, ink 304, not transcribed) ends on
% « rendre compte. de ces des singularités qu'ils présente » and view 622
% (page 2) continues it with « J'exposerai ensuite celles de mes
% experiences »; view 621 between them is blank. Batch 31 reads v. 620 the
% same way (ink 304, the draft stopping on « qu'ils présente », v. 621
% blank) and takes v. 614–620 for drafts of the opening of a prize memoir
% addressed to « la classe »; v. 622–626 are therefore that draft's pages
% 2–4. Beside 304 it reads « a small blotted figure, perhaps 1, struck and
% underlined »: this pass reads the same mark as her page number 1,
% struck, since 2, 3, 4 follow on v. 622–626. 640/641: view 641 is the
% blank back of v. 640; v. 640 ends on an unfinished list of areas. Batch
% 33 opens at v. 643 (a slip inked 315) with a change of variables for v
% about an inclined axis, not with the sequel of v. 640.
%
% The numbers on the leaves, and why no \folio{} is written. (1) The ink
% series at the top right, in the larger rounder hand: 305 (v. 622), 306
% (624), 307 (626), 308 (628, a blank leaf), 309 (630), 310 (632), 311
% (634, the blank side of the leaf whose other side is v. 635), 312 (636),
% 313 (638, the blank side of the leaf whose other side is v. 639), 314
% (640); 304 on v. 620 outside the batch. One per physical leaf, blank
% leaf included, without a gap; the first figure on v. 622 has a
% descending tail and is read 5 by the run. On v. 639 a « 312 » struck
% through with a horizontal stroke stands on the written side of the leaf
% inked 313 on its blank side: that leaf was numbered twice. Penned, not
% pencilled, as everywhere in this volume: no \folio{}. (2) Her own page
% numbers, at the head: 2 (v. 622, underlined and struck obliquely), 3
% (624, underlined), 4 (626, underlined and struck), then on the heat
% leaves 7 (632, struck), 8 (635, not struck, no underline), 7 (636,
% struck), 6 (639, struck). Two leaves read 7; recorded as read. The « 8 »
% at the top right of v. 637 and 639 stands on a narrow strip of paper
% showing above those leaves and seems to be the head of the page of
% v. 635; it is not a number of those pages. (3) Numbers in the text:
% « N.o » without a figure at v. 622 and 624; « N.o 11 de mon memoire pour
% 1813 » (v. 626). Every number above was read on its view; none is
% inferred. None of these leaves can be Laubenbacher and Pengelley's
% « Manuscript C » (ff. 348r–349r): the subject is elastic surfaces and
% heat, and the ink series runs 305–314.
%
% Notation. Hers, kept. « sin », « cos » (with or without a stop) are set
% \text{sin}, \text{cos}; the stop after them is not recorded stop by
% stop, but the stop of « d. » before a function (d.sin θ, d.(r² dv/dr))
% is kept, as is the stop she uses for multiplication. Her π is a tall Π
% with unequal legs that reads almost « 17 » (v. 632); it is set \pi. Her
% exponent 2 and exponent 4 are close (the 4 written like « b »); where
% the calculation and the glyph pull apart the reading is marked
% \uncertain{}. The letter read D in K/C.D is a dome on a bar. Words she
% underlines in the prose are set in \emph{}. Her « Retournez » at the foot
% of the margin of v. 636 is kept. The small flower drawings in the margin
% of v. 632 are described in a note, not reproduced.
% Spelling is hers and is not flagged: « quarrées », « correspondans »,
% « connoître », « parallels », « étoit », « devoit », « tems »,
% « parconséquent », « apartenir », « equivallante », « c'est-adire »,
% « dela », « debase », « alaquelle », « aumoins », « cequi », « lemême »,
% « leplan », « ontrouve », « jusqu'àprésent », « en fin », « vis avis »,
% « deux ligne droites », « deux lignes droits », « aux des côtés ».
% Where the leaf and the calculation disagree the leaf stands, with a note
% (v. 630: z = cos θ without r; v. 639: dv/dr for dv/dt, a missing dφ).
%
% Deletions. Horizontal strokes through words are given as \struck{}. Two
% passages of v. 622 are cancelled by large oblique strokes, and three
% lines of v. 626 by crosses over each word: the first are said to be so
% in a note, the second are wrapped in \struck{} with a note.
%
% What could not be read, and what is offered with doubt. \ill{} stands
% four times: a group scribbled over at the end of a line of d²v/dx²
% (v. 630, struck); the variable after « la seconde demême par rapport a »
% (v. 636), which the sense wants to be sin θ; a blotted numerator (v. 639)
% and a blotted factor (v. 640). \uncertain{} stands for « ployés »,
% « cambrer », a struck « et », « sauf » (v. 622); « est », the 2 in the
% margin addition, the exponents 4 and 2 (v. 630); y, z, « cone », « fera »
% (v. 632); « surbaissé » and the letter r (v. 635); « obliquement », « et
% sort » (v. 636); « ce » (v. 639); v in dv/dθ (v. 640). Nothing was
% guessed to fill a gap.
%
% Nothing here was checked against any published edition of Sophie
% Germain's papers, of Chladni, or of any memoir on heat.
\documentclass[11pt,a4paper]{article}
% The preamble shared by every transcription of the mathematicians' manuscripts in Gallica.
%
% It rests on one idea: **uncertainty compounds, it does not resolve.** A
% transcription that smooths over a doubtful word has destroyed the only thing
% it was made for — a reader must be able to tell, at every point, what was on
% the page and what was supplied. The apparatus macros below are therefore the
% critical apparatus, not decoration.
%
% The same commands are recognised by `scripts/render.mjs`, which produces the
% site's reading view, and by `scripts/tei.mjs`. A macro added here must be
% added there too, or rendering fails loudly — which is the wanted behaviour.
%
% Comments here are English, like the rest of the repository. The documents
% this preamble sets are French, and that is deliberate: see the skills.

\ProvidesPackage{germain}[2026/09/15 Transcriptions of mathematicians' manuscripts in Gallica]

\usepackage{amsmath,amssymb,amsthm}
\usepackage{mathrsfs}
% Kept from the parent preamble so the renderer's diagram support stays
% available; these manuscripts draw no commutative diagrams, and nothing here uses it.
\usepackage{tikz-cd}
\usepackage[english,french]{babel}
\usepackage{fontspec}

% --- Greek ------------------------------------------------------------------
% The text face stays fontspec's default, Latin Modern, which has no Greek:
% XeTeX drops every glyph it lacks with a line in the log and nothing on the
% page, and the Greek words the manuscripts quote vanished from the PDFs. So
% the Greek and Greek Extended blocks get a XeTeX character class of their own,
% and entering or leaving a run of it switches the family to CMU Serif — the
% same Computer Modern design, with polytonic Greek — and back. Only the
% family changes: a Greek word in \textit or \textbf keeps its shape and series.
% The transitions to and from every other class carry the tokens that class
% has towards an ordinary letter, so French spacing before « ; » or inside
% « » still holds after a Greek word. No grouping, so no brace or \endgroup
% can come out unbalanced; maths is left alone, since XeTeX inserts nothing
% there. Kept identical in `exercices.sty`.
\makeatletter
\newfontfamily\germainGreekfont{cmun}[
  Extension=.otf, NFSSFamily=germaingreek,
  UprightFont=*rm, ItalicFont=*ti, SlantedFont=*sl,
  BoldFont=*bx, BoldItalicFont=*bi, BoldSlantedFont=*bl]
\newXeTeXintercharclass\germain@greek
\def\germain@greekrange#1#2{%
  \@tempcnta=#1\relax
  \loop
    \XeTeXcharclass\@tempcnta=\germain@greek
    \ifnum\@tempcnta<#2\advance\@tempcnta\@ne\repeat}
\germain@greekrange{"0370}{"03FF}
\germain@greekrange{"1F00}{"1FFF}
\def\germain@greekprev{\rmdefault}
\def\germain@greekin{%
  \edef\germain@greekprev{\f@family}\fontfamily{germaingreek}\selectfont}
\def\germain@greekout{\fontfamily{\germain@greekprev}\selectfont}
\def\germain@greekjoin{%
  \XeTeXinterchartokenstate=\@ne
  \@tempcnta=\z@
  \loop
    \ifnum\@tempcnta=\germain@greek\else
      \XeTeXinterchartoks\germain@greek\@tempcnta=\expandafter{%
        \expandafter\germain@greekout\the\XeTeXinterchartoks\z@\@tempcnta}%
      \XeTeXinterchartoks\@tempcnta\germain@greek=\expandafter{%
        \the\XeTeXinterchartoks\@tempcnta\z@\germain@greekin}%
    \fi
    \ifnum\@tempcnta<4095\advance\@tempcnta\@ne\repeat}
% babel-french sets its own classes, and saves and restores the interchar
% state, each time a language is selected: join again after it does.
\addto\extrasfrench{\germain@greekjoin}
\addto\extrasenglish{\germain@greekjoin}
\AtBeginDocument{\germain@greekjoin}
\makeatother

\usepackage{geometry}
\geometry{a4paper,margin=25mm,marginparwidth=28mm}
\usepackage{marginnote}
\usepackage{xcolor}
\usepackage[normalem]{ulem}
\setlength{\emergencystretch}{3em}
\usepackage{hyperref}
\hypersetup{colorlinks=true,linkcolor=black,urlcolor=black,pdfborder={0 0 0}}

% --- Batch metadata ---------------------------------------------------------
% Taken as they stand by the rendering script, which shows them at the head of
% the reading view. They fix what the file claims to cover. The macro names are
% the parent project's, so that the scripts stay shared; \folder here means the
% volume's slug — `fr-9115` — and \pages counts Gallica views.

\newcommand{\folder}[1]{\gdef\grothFolder{#1}}
\newcommand{\batch}[1]{\gdef\grothBatch{#1}}
\newcommand{\pages}[2]{\gdef\grothFirst{#1}\gdef\grothLast{#2}}
\newcommand{\foldertitle}[1]{\gdef\grothFoldertitle{#1}}

% The holder's own shelfmark, printed as they write it: « Français 9115 ».
\newcommand{\shelfmark}[1]{\gdef\grothShelfmark{#1}}

% The dating, copied verbatim from the holder's catalogue — never inferred
% here. « XIXe siècle » is the BnF's; a pass that dates a leaf from its content
% says so in a \note{}, as its own inference.
\newcommand{\dating}[1]{\gdef\grothDating{#1}}

% The status notice, printed in the title block and repeated as a diagonal
% watermark on every page of the PDF. The manuscripts are in the public
% domain; what this line declares is not a right but a quality — an
% unverified first pass by a machine. The skills write the line; a file
% without it gets no watermark, so leaving it out is visible in the output.
\newcommand{\watermark}[1]{\gdef\grothWatermark{#1}}

\gdef\grothFolder{?}\gdef\grothBatch{}\gdef\grothFirst{?}\gdef\grothLast{?}%
\gdef\grothFoldertitle{}\gdef\grothShelfmark{}\gdef\grothDating{}\gdef\grothWatermark{}

% --- The résumé -------------------------------------------------------------
\newenvironment{resume}
  {\small\setlength{\parskip}{0.45em}}
  {\par\medskip\hrule\medskip}

% --- Keywords ---------------------------------------------------------------
% In English, closing the modernised reading's résumé: the modern vocabulary
% under which to search for what the leaves construct. The manifest extracts
% them to tag the volume — this line is the single source of the tags.
\newcommand{\keywords}[1]{\par\smallskip\noindent
  {\footnotesize\textsc{Keywords} --- \textit{#1}}\par}

% --- The critical apparatus -------------------------------------------------

% The Gallica view number, in the margin. This is the anchor that lets a
% disagreement over a formula be settled by looking at *one* image rather than
% a volume of 750. The site uses it to turn the facsimile as the transcription
% is scrolled. A view is one scanned image: usually one side of a leaf.
\newcommand{\page}[1]{%
  \marginnote{\footnotesize\color{gray!70}\textsf{v.\,#1}}%
  \label{page:#1}\ignorespaces}

% The BnF's pencilled foliation, where the leaf carries it and a pass can read
% it: « 198r ». This is what the literature cites — « ff. 198r–208v » — and
% the one line that turns a citation into a view. Recorded, never inferred:
% a folio a pass cannot read is not written.
\newcommand{\folio}[1]{%
  \marginnote{\footnotesize\color{gray!50}\textsf{f.\,#1}}\ignorespaces}

% The modernised reading's coarser counterpart to \page{}: which views a
% section is drawn from.
\newcommand{\pagerange}[2]{%
  \marginnote{\footnotesize\color{gray!70}\textsf{v.\,#1--#2}}%
  \ignorespaces}

% Illegible: never guessed.
\newcommand{\ill}{\textcolor{red!60!black}{[\,\ldots\,]}}

% A doubtful reading: the word is given, and flagged as uncertain.
\newcommand{\uncertain}[1]{\uline{#1}}

% An editorial addition — a missing letter, an implied word.
\newcommand{\add}[1]{\textcolor{blue!50!black}{[#1]}}

% Struck out by the writer. What was crossed out is part of the document.
\newcommand{\struck}[1]{\sout{#1}}

% The transcriber's note, on the state of the page or the mathematics.
\newcommand{\note}[1]{\par\smallskip{\small\color{gray!60!black}\textit{[#1]}}\par\smallskip}

% A marginal note of hers, distinct from the body of the text.
\newcommand{\marginal}[1]{\par\smallskip\noindent\hspace{1em}%
  {\small\color{gray!60!black}#1}\par\smallskip}

% --- Title ------------------------------------------------------------------
% The title block is French, like the documents themselves.

\AddToHook{shipout/background}{%
  \ifx\grothWatermark\empty\else
    \begin{tikzpicture}[remember picture, overlay]
      \node[rotate=52, scale=3.4, align=center, text=gray!18,
            font=\sffamily\bfseries]
        at (current page.center) {\grothWatermark};
    \end{tikzpicture}%
  \fi}

\AtBeginDocument{%
  \begin{center}
    {\large\bfseries Manuscrits de mathématiciens (Gallica) ---
      \ifx\grothShelfmark\empty\grothFolder\else\grothShelfmark\fi}\\[2pt]
    {\itshape\grothFoldertitle}\\[4pt]
    {\small vues \grothFirst--\grothLast
      \ifx\grothBatch\empty\else\ (lot \grothBatch)\fi}\\[2pt]
    \ifx\grothDating\empty\else
      {\small\color{gray!60!black}Datation du catalogue : \grothDating}\\[2pt]
    \fi
    \ifx\grothWatermark\empty\else
      {\small\bfseries\color{red!45!gray}\grothWatermark}\\[2pt]
    \fi
    {\footnotesize\color{gray!60!black}Transcription établie d'après les images IIIF de Gallica.
     Source gallica.bnf.fr / Bibliothèque nationale de France.\\
     Les lectures incertaines sont soulignées, les illisibles notées [\,\ldots\,],
     les ajouts entre crochets ; \textsf{v.} numérote les vues de Gallica,
     \textsf{f.} la foliotation au crayon de la BnF.}
  \end{center}
  \vspace{1em}}

\endinput


\folder{fr-9115}
\shelfmark{Français 9115}
\batch{32}
\pages{621}{640}
\dating{1801-1900}
\watermark{Lecture automatique\\non vérifiée}
\foldertitle{Papiers de Sophie GERMAIN. — Recueil de dissertations et problèmes mathématiques et physiques.}

\begin{document}

\note{Numérotation des feuillets. En haut à droite, à l'encre, d'une écriture plus ronde, la série qui court dans tout le volume : 305 (vue 622), 306 (624), 307 (626), 308 (628, feuillet blanc), 309 (630), 310 (632), 311 (634, face blanche du feuillet dont l'autre face est la vue 635), 312 (636), 313 (638, face blanche du feuillet dont l'autre face est la vue 639), 314 (640). Sur la vue 639, un « 312 » barré d'un trait horizontal, sur la face écrite du feuillet numéroté 313 sur l'autre face. Cette série numérote des feuillets, feuillet blanc compris ; elle occupe la place de la foliotation de la Bibliothèque, mais elle est tracée à la plume et non au crayon, et aucun « f. » n'est donc porté. Sa pagination, en tête : 2 (vue 622), 3 (624), 4 (626), puis 7 (632), 8 (635), 7 (636), 6 (639). Tous ces nombres ont été lus sur la vue ; aucun n'est déduit. Ne sont pas transcrites : les vues 621, 623, 625, 627, 631 et 633, dos blancs des feuillets qui les précèdent, où l'écriture se lit à l'envers par transparence ; les vues 628 et 629, feuillet 308 blanc des deux côtés ; les vues 634 et 638, faces des feuillets 311 et 313 qui ne portent que ce numéro.}

\page{622}

\struck{je donnerai ensuite les} J'exposerai ensuite celles de mes experiences
dont j'ai tenu note en fin je \struck{monte} comparerai le petit
nombre defaits qui s'y prêtent aux resultats du calcul
dans le cas le plus simple qui puisse se présenter
\note{En haut à droite, à l'encre, « 305 », le dernier chiffre à queue descendante (lu 5 parce que le feuillet précédent, vue 620, hors de ce lot, porte 304) ; à côté, souligné et barré d'un trait oblique, « 2 », sa pagination. La phrase continue la page 1 (vue 620), qui s'arrête sur « rendre compte. de ces des singularités qu'ils présente ». « le petit nombre » : le « t » final de « petit » porte le paraphe qu'elle met à la fin des mots, lu sans « s ».}

N.o Dans la vue de rendre, \struck{autant que possible les faits}
\add{\struck{entre} les experiences que j'allois \struck{entreprendre et celles qui}}
experiences sur les surfaces cylindriques comparables \struck{a celles}
\add{au cas le plus simple que \struck{presente} la vibration des \struck{plaques} surfaces planes}
\struck{qui ont été} faits \struck{sur les} plaques \struck{quarrées}
\add{\struck{ont été faits sur les surfac.}}, je me \struck{suis} procuré
des portions cylindriques dont la hauteur avoit pour mesure
l'étendue de l'arc A et dont les quatre côtés étoient
deux a deux exactement parallels.
\note{Le « N.o » est en tête du paragraphe, sans chiffre. Les deux interlignes sont écrits au-dessus des lignes 2 et 3 et au-dessous de la ligne 3 ; on les donne à la place où l'œil les prend. La phrase, dans l'état où la laissent les corrections, n'est pas achevée : « Dans la vue de rendre les experiences sur les surfaces cylindriques comparables au cas le plus simple [que] la vibration des surfaces planes ». À partir de « je me suis procuré » jusqu'à « parallels. », le passage est annulé par de grands traits obliques croisés, et non par des traits horizontaux.}

En \struck{tirant la pla surface} appliquant les doigts au centre
de figure de la surface et tirant le son d'un des angles
on obtient \struck{que g} toujours, quelque soit le nombre de
degrés compris dans l'arc A, une figure nodale composée
\note{Ce paragraphe entier est annulé par quatre grands traits obliques qui le traversent de haut en bas ; la phrase s'arrête sur « composée » en fin de ligne, sans suite.}

J'ai fait taillé un certain nombre de plaques quarrées
qui ont été ensuite \uncertain{ployés} suivant \add{les} differentes courbures
dont je voulois \struck{\uncertain{et}} connoître l'influence.
\note{« ployés » : la fin du mot est récrite et peut se lire « ployées ». « taillé » : ainsi, avec l'accent.}

Avant de faire \uncertain{cambrer} les plaques dont il s'agit
j'avois essayé \struck{tous} les tons correspondans aux \struck{diff deux}
figures \struck{les plus simples qui se forment dans le}
nodales qui peuvent s'y former j'esperois obtenir par là
une première indication de l'effet dela courbure
\note{« cambrer » : le mot commence par une lettre récrite ; on lit « camber » ou « cambrer ».}

j'ai vu avec étonnement que la première figure \add{nodale} de
Mr Chladni c'est-adire celle qui est \struck{formée} \add{composée} de
deux ligne droites qui se croisent au centre de
la plaque figure qui se forme \struck{a} de la même
manière sur les surface cylindrique, que sur les \struck{surf}
plaque étoit \struck{\uncertain{sauf} de très petites} accompagnée du
\note{« ligne », « surface », « plaque » : sans « s » sur le feuillet. Le premier mot barré de la dernière ligne a le s long ; « de très » est écrit en un mot. La page s'arrête sur « accompagnée du » ; la suite est à la vue 624.}

\page{624}

même ton. après et avant la courbure, sur une quarantaine
de surfaces qui ont été ainsi essayées il s'est trouvé
sept ou huit fois une difference d'un demi ton, le
son de la surface courbée étoit alors le plus élevé
Je suis porté a croire que le contraire devoit arriver
je pense que la haute temperature alaquelle on est
forcé d'elever le verre pour le rendre flexible complique
l'effet que je voulois obtenir en changeant l'élasticité
naturelle des pièces. Le peu d'activité des ateliers
a l'époque ou je me suis occupé de ce genre
\struck{d'expe} d'essais m'a empêché de les varier sous \add{les} differentes
conditions que j'avois imaginées pour en écarter toute
cause d'erreurs.
\note{En tête, au milieu, « 3 » souligné, sa pagination ; en haut à droite, à l'encre, « 306 ». Le cachet de la Bibliothèque impériale est frappé au milieu du feuillet. La première ligne continue la page 2 (vue 622), qui s'arrête sur « accompagnée du » : « accompagnée du même ton ».}

N.o En essayant de suite aumoins quarante plaques
quarrées j'ai eu occasion de m'assurer qu'un des
intervalle le plus constant est celui des sons qui
accompagnent sur ce genre desurface les figures
nodales composées \struck{d'une ligne tracé} de deux
lignes droites perpendiculaires l'une a l'autre
et de \struck{deux} trois lignes egalement droites dont
l'une est perpendiculaire aux deux autres
je n'ai remarqué que deux fois \struck{une} difference
qui n'étoit pas même d'un demi ton. *
\marginal{Lorsqu'il s'agit des surfaces cylindriques l'intervalle des sons qui accompagne les figures analogues a celles dont je viens de parler \struck{est touj} varie en raison du plus ou moins de courbure.}
\note{La note marginale est écrite dans la marge de gauche, en face des lignes « nodales composées » à « qui n'étoit pas même d'un demi ton » ; aucun signe ne l'appelle dans la marge, mais l'astérisque qui suit « demi ton. » paraît la désigner. « N.o » sans chiffre, comme à la page précédente. « un des intervalle le plus constant » : ainsi, au singulier.}

\struck{On sait que} Sur les plaques quarrées \struck{le} un
\struck{même} \add{donné} nombre de lignes nodales est toujours
accompagné \struck{toujours} du même ton, \struck{il se}
\struck{fait même} ce principe est si constant
\note{« ce » est récrit sur un autre mot, sans doute « le ». La page s'arrête sur « si constant », au tiers inférieur du feuillet, le reste blanc.}

\page{626}

qu'il a guidé Mr Chladni dans la classification
\struck{ingénieuse qu'il a faite} des distorsions des figures
nodales et la \struck{Th} théorie a pleinement \add{confirmé}, comme
je l'ai dit ailleurs, les vues ingénieuses de ce
physicien.
\note{En tête, au milieu, « 4 » souligné et barré d'un trait oblique ; en haut à droite, à l'encre, « 307 ». La première ligne continue la page 3 (vue 624), qui s'arrête sur « ce principe est si constant ».}

Le contraire arrive dans les surfaces cylindriques
\struck{le même on verra par le détail des résultats faits que}
\struck{dont je donnerai plus bas la note} \struck{j'exposerai plus bas}
\struck{l'exposition des fa} On verra toute a l'heure que
\struck{les lignes paralleles au côté courbe dela surface}
\struck{ne sont pas équi} que des figures \add{nodales} composées
du même nombre de lignes paralleles \add{les unes aux côtés courbes} \struck{au côté}
\add{les autres aux côtés droit dela même surface} de la surface \struck{et} sont accompagnées de
sons differens.
\note{Les deux interlignes « les unes aux côtés courbes » et « les autres aux côtés droit dela même surface » sont écrits au-dessus des deux dernières lignes du paragraphe ; « de la surface » reste sur la ligne, non barré.}

En fin la figure \add{nodale} \struck{com} composée de deux lignes diagonales
\struck{qui} se montre \struck{très} \add{très} facilement sur les plaques quarrées
\struck{ne peut se former sur les surface cylindriques}
\struck{il suffit d'une courbure a peine sensible pour}
\struck{que cette figure ne paroisse plus.} \struck{Selon} suivant
\struck{l'indication de Mr Chladni et} la théorie que \add{\struck{dont il s'agit}}
j'ai developpée N.o 11 de mon memoire pour 1813, \struck{cette}
la figure \add{dont il s'agit} est equivallante a celle qui seroit composée
de deux lignes droits paralleles aux des côtés dela
plaque. Mr Chladni a remarqué que jamais
cette dernière figure \struck{qui} ne se forme sur les plaques
quarrées mais qu'elle est toujours remplacée par
les distorsions dont \struck{j'ai expliqué ailleurs} la nature
a été expliquée ailleurs
\note{Les trois lignes « ne peut se former … ne paroisse plus. » sont barrées mot à mot de croix obliques, non d'un trait horizontal. « très » est barré sur la ligne et récrit au-dessus. Le premier « dont il s'agit », au-dessus de « 1813, cette » à droite, est barré ; il est récrit au-dessus de « figure est ». « la », qui remplace « cette » barré, est écrit dans la marge de gauche, devant « figure ». Le « mémoire pour 1813 » et son N.o 11 sont nommés ainsi sur le feuillet ; on n'identifie pas ici ce mémoire. « deux lignes droits » et « aux des côtés » : ainsi. La page s'arrête sur « a été expliquée ailleurs », au bas du feuillet.}

\page{630}

\emph{cône}

on prendra, comme lorsqu'il s'est agi dela sphère, ou cequi \uncertain{est} la même
chose dela \struck{surfa} du cone pris dans la sphère et terminé non pas
par une base plane mais la calotte sphérique qui a sa base commune
avec celle du cône.
\note{En haut à droite, à l'encre, « 309 » ; aucune pagination de sa main. Le titre « cône » est souligné, en tête et au milieu. Le mot lu « est » après « cequi » est empâté. La phrase n'a pas de complément à « on prendra » : ainsi sur le feuillet. Le feuillet ne continue pas la page 4 (vue 626) : c'est un autre papier, sans numéro de page, et un autre sujet.}

\[ x = r\,\text{sin}\,\theta\,\text{cos}\,\varphi, \qquad y = r\,\text{sin}\,\theta\,\text{sin}\,\varphi, \qquad z = \text{cos}\,\theta. \]
\[ dx = dr\,\text{sin}\,\theta\,\text{cos}\,\varphi + r\,\text{cos}\,\theta\,\text{cos}\,\varphi\,d\theta - r\,\text{sin}\,\theta\,\text{sin}\,\varphi\,d\varphi. \]
\[ dy = dr\,\text{sin}\,\theta\,\text{sin}\,\varphi + r\,\text{cos}\,\theta\,\text{sin}\,\varphi\,d\theta + r\,\text{sin}\,\theta\,\text{cos}\,\varphi\,d\varphi \]
\[ dz = dr\,\text{cos}\,\theta - r\,\text{sin}\,\theta\,d\theta. \]
\note{« $z = \text{cos}\,\theta$ » : sans le facteur $r$, que la différentielle $dz$ de la ligne suivante suppose ; ainsi sur le feuillet.}

avant de procéder aux transformations qui doivent donner les valeurs de $\frac{d^2v}{dx^2} + \frac{d^2v}{dy^2}$
$+ \frac{d^2v}{dz^2}$ relatives au cône il faut observer que les differentielles ne doivent pas
être prise ici par rapport a l'arc $\theta$, mais par rapport au sinus
de cet angle; la raison de ce changement est que l'arc $\theta$
qui mesure la calotte sphérique est remplacée ici par $\text{sin}\,\theta$ qui mesure
la base du cône qui considéré dans la sphère s'appuyoit sur
la calotte sphérique.
\note{« les » est récrit sur un autre mot. Le dernier dénominateur, au début de la deuxième ligne, est tracé presque comme « $dx^2$ » ; on le lit $dz^2$.}

cela posé on aura
\[ \frac{dr}{dx} = \text{sin}\,\theta\,\text{cos}\,\varphi, \quad \frac{d.\text{sin}\,\theta}{dx} = \frac{\text{cos}^2\theta\,\text{cos}\,\varphi}{r}, \quad \frac{d\varphi}{dx} = -\frac{\text{sin}\,\varphi}{r\,\text{sin}\,\theta} \]
\[ \frac{dr}{dy} = \text{sin}\,\theta\,\text{sin}\,\varphi, \quad \frac{d.\text{sin}\,\theta}{dy} = \frac{\text{cos}^2\theta\,\text{sin}\,\varphi}{r}, \quad \frac{d\varphi}{dy} = \frac{\text{cos}\,\varphi}{r\,\text{sin}\,\theta} \]
\[ \frac{dr}{dz} = \text{cos}\,\theta, \quad \frac{d.\text{sin}\,\theta}{dz} = -\frac{\text{sin}\,\theta\,\text{cos}\,\theta}{r}, \]

\[ \frac{dv}{dx} = \frac{dv}{dr}.\text{sin}\,\theta\,\text{cos}\,\varphi + \frac{dv}{d.\text{sin}\,\theta}.\frac{\text{cos}^2\theta\,\text{cos}\,\varphi}{r} - \frac{dv}{d\varphi}.\frac{\text{sin}\,\varphi}{r\,\text{sin}\,\theta} \]
\[ \frac{dv}{dy} = \frac{dv}{dr}.\text{sin}\,\theta\,\text{sin}\,\varphi + \frac{dv}{d.\text{sin}\,\theta}.\frac{\text{cos}^2\theta\,\text{sin}\,\varphi}{r} + \frac{dv}{d\varphi}.\frac{\text{cos}\,\varphi}{r\,\text{sin}\,\theta} \]
\[ \frac{dv}{dz} = \frac{dv}{dr}.\text{cos}\,\theta - \frac{dv}{d.\text{sin}\,\theta}.\frac{\text{cos}\,\theta\,\text{sin}\,\theta}{r} \]
\note{Dans $\frac{dv}{dy}$, le « cos » du dernier numérateur est repassé et empâté. Le cachet de la Bibliothèque impériale (« MSS. ») est frappé à droite de ces trois lignes.}

et en mettant $\sqrt{1-\text{sin}^2\theta}$ a la place de $\text{cos}\,\theta$ \quad $1-\text{sin}^2\theta$ a la place de $\text{cos}^2\theta$.

\begin{align*}
\frac{d^2v}{dx^2} = {} & \text{sin}\,\theta\,\text{cos}\,\varphi\left\{\frac{d^2v}{dr^2}.\frac{dr}{dx} + \frac{d^2v}{dr\,d.\text{sin}\,\theta}.\frac{d.\text{sin}\,\theta}{dx} + \frac{d^2v}{dr\,d\varphi}.\frac{d\varphi}{dx}\right\} \\
& + \frac{\text{cos}^2\theta\,\text{cos}\,\varphi}{r}\left\{\frac{d^2v}{d.\text{sin}\,\theta.dr}\frac{dr}{dx} + \frac{d^2v}{d.\text{sin}\,\theta.d\varphi}.\frac{d\varphi}{dx} + \frac{d^2v}{(d.\text{sin}\,\theta)^2}.\frac{d.\text{sin}\,\theta}{dx}\right\} \\
& - \frac{\text{sin}\,\varphi}{r\,\text{sin}\,\theta}\left\{\frac{d^2v}{d\varphi.dr}.\frac{dr}{dx} + \frac{d^2v}{d\varphi.d.\text{sin}\,\theta}.\frac{d.\text{sin}\,\theta}{dx} + \frac{d^2v}{d\varphi^2}.\frac{d\varphi}{dx}\right\} \\
& + \frac{dv}{dr}\left\{\frac{d.\text{sin}\,\theta}{dx}.\text{cos}\,\varphi - \frac{d\varphi}{dx}.\text{sin}\,\theta\,\text{sin}\,\varphi\right\} \\
& - \frac{dv}{d.\text{sin}\,\theta}\left\{\frac{d\varphi}{dx}.\frac{\text{sin}\,\varphi\,\text{cos}^2\theta}{r} + \frac{dr}{dx}\,\frac{\text{cos}^2\theta\,\text{cos}\,\varphi}{r^2}\ \struck{\ill{}}\ \add{-\frac{d.\text{sin}\,\theta}{dx}\,\frac{\uncertain{2}\,\text{sin}\,\theta\,\text{cos}\,\varphi}{r}}\right\} \\
& - \frac{dv}{d\varphi}\left\{\frac{d\varphi}{dx}.\frac{\text{cos}\,\varphi}{r\,\text{sin}\,\theta} - \frac{d.\text{sin}\,\theta}{dx}.\frac{\text{sin}\,\varphi}{r\,\text{sin}^2\theta} - \frac{dr}{dx}.\frac{\text{sin}\,\varphi}{r^2\,\text{sin}\,\theta}\right\}
\end{align*}
\note{Première ligne : le $\varphi$ du dernier $\frac{d\varphi}{dx}$ est écrit sur un autre signe et empâté. Deuxième ligne : « $\frac{dr}{dx}$ » est serré après le premier dénominateur et le « $+$ » qui suit est empâté. Troisième ligne : le $r$ du premier $\frac{dr}{dx}$ est empâté. Quatrième ligne : entre « $\text{sin}\,\theta$ » et « $\varphi$ » un « sin » récrit. Cinquième ligne : à la fin, un groupe écrit puis noirci de boucles, illisible, au-dessous d'un ajout écrit au-dessus de la ligne jusqu'au bord droit et fermé d'un trait vertical, « $-\frac{d.\text{sin}\,\theta}{dx}$ » suivi d'une fraction dont le numérateur, lu « $2\,\text{sin}\,\theta$ », a été récrit (un « cos » barré à la fin) et dont la seconde ligne porte « $\text{cos}\,\varphi$ » ; le chiffre lu 2 peut être un 3.}

\begin{align*}
= {} & \text{sin}^2\theta\,\text{cos}^2\varphi.\frac{d^2v}{dr^2} + \frac{2\,\text{cos}^2\theta\,\text{sin}\,\theta\,\text{cos}^2\varphi}{r}.\frac{d^2v}{dr\,d.\text{sin}\,\theta} - \frac{2\,\text{sin}\,\varphi\,\text{cos}\,\varphi}{r}\,\frac{d^2v}{dr\,d\varphi} \\
& + \frac{\text{cos}^{\uncertain{4}}\theta\,\text{cos}^2\varphi}{r^2}.\frac{d^2v}{(d.\text{sin}\,\theta)^2} - \frac{2\,\text{cos}^2\theta\,\text{cos}\,\varphi\,\text{sin}\,\varphi}{r^2\,\text{sin}\,\theta}.\frac{d^2v}{d\varphi\,d.\text{sin}\,\theta} + \frac{\text{sin}^2\varphi}{r^2\,\text{sin}^2\theta}.\frac{d^2v}{d\varphi^2} \\
& + \frac{dv}{dr}\left\{\frac{\text{cos}^2\theta\,\text{cos}^2\varphi}{r} + \frac{\text{sin}^2\varphi}{r}\right\} \\
& + \frac{dv}{d.\text{sin}\,\theta}\left\{\frac{\text{sin}^2\varphi\,\text{cos}^2\theta}{r^2\,\text{sin}\,\theta} - \frac{\text{sin}\,\theta\,\text{cos}^2\theta\,\text{cos}^2\varphi}{r^2} - \frac{2\,\text{sin}\,\theta\,\text{cos}^{\uncertain{2}}\theta\,\text{cos}^2\varphi}{r^2}\right\} \\
& + \frac{dv}{d\varphi}\left\{\frac{\text{cos}\,\varphi\,\text{sin}\,\varphi}{r^2\,\text{sin}^2\theta} + \frac{\text{cos}^2\theta\,\text{sin}\,\varphi\,\text{cos}\,\varphi}{r^2\,\text{sin}^2\theta} + \frac{\text{sin}\,\varphi\,\text{cos}\,\varphi}{r^2}\right\}
\end{align*}
\note{Le « sin » qui ouvre la ligne est récrit. Ses exposants 2 et 4 ont presque la même forme (le 4 tracé comme un « b ») : l'exposant de $\text{cos}\,\theta$ dans le quatrième terme est lu 4, celui du dernier terme de l'accolade de $\frac{dv}{d.\text{sin}\,\theta}$ est récrit et reste douteux. La page s'arrête sur l'accolade de $\frac{dv}{d\varphi}$, au bas du feuillet ; le calcul de $\frac{d^2v}{dy^2}$ n'est pas sur cette page.}

\page{632}

Si on mène par le \struck{autre} sommet du cône opposé un plan
perpendiculaire a une position donné du rayon $r$ qui est
le côté de ce cone, lorsqu'il s'agit de $r = z$ ce plan est
celui des $x, y$ s'il s'agit ce plan
\note{En haut à droite, à l'encre, « 310 » ; à sa gauche, « 7 » traversé d'un trait oblique, sans doute sa pagination barrée, comme le « 2 » de la vue 622. La phrase s'arrête sur « s'il s'agit ce plan », sans suite ; un trait horizontal court la sépare de ce qui suit. Dans la marge de gauche, sur toute la hauteur du texte, de petits dessins à la plume : des fleurs, une feuille, un rameau feuillu ; on les signale sans les décrire un à un.}

Soit $A$ comme jusqu'àprésent l'angle que le côté du
cône donné fait avec l'axe et $A'$ l'angle que le côté
du cône opposé fait avec lemême axe. le plan perpendiculaire
a une position donnée du côté \add{$r = z$} du cône donné
coupera le plan des $zx$ et celui des $xy$. \struck{et} l'angle que
l'intersection avec le plan des $zx$ fait avec l'axe des $x$ est
$\frac{A\,v}{\frac{\pi}{2}}$. et celui que l'intersection du même plan perpendiculaire
au rayon $r$ avec \struck{l'axe} le plan de $\uncertain{y}x$ fait avec l'axe
des $x$ est $\struck{\frac{A\,v}{\frac{\pi}{2}}}$. $\frac{\pi}{2} - \frac{A\,v}{\frac{\pi}{2}}$
\marginal{N.ota on a $A + A' = \frac{\pi}{2}$}
\note{La note marginale est écrite à gauche, en face des lignes « du cône opposé … coupera le plan », après une grosse tache d'encre cernée d'un trait. Son $\pi$ est tracé comme un grand $\Pi$ aux jambages inégaux, qui se lit presque « 17 » ; on le rend partout par $\pi$. Le numérateur lu « $A\,v$ » : un $A$ suivi d'une lettre bouclée surmontée d'un petit trait, lue $v$ ; ce que désigne $v$ n'est pas dit sur la page. « $r = z$ » est écrit au-dessous de « du côté ». « avec » après « rayon $r$ » est repassé ; le $y$ de « le plan de $yx$ » est empâté, et le long trait oblique qui en part descend barrer la première fraction de la ligne suivante.}

l'angle que l'intersection du plan perpendiculaire
a une position donné de $r = \uncertain{z}$ \struck{du} qui est le côté
du \uncertain{cone} opposé \struck{fait sur le pl} sur le plan de \struck{$yx$}
$zx$ fait avec \struck{l'angle des} l'axe des $x$ est $\frac{\pi}{2} - \frac{A\,v}{\frac{\pi}{2}}$
et l'angle que l'intersection du même plan \uncertain{fera}
celui de $xy$ fait avec l'axe de $x$ est $\frac{A\,v}{\frac{\pi}{2}}$.
\note{« $r = z$ » : la lettre après le signe est récrite et paraît barrée ; « cone » est suivi d'une syllabe qu'on ne sépare pas du mot (« conedu ») ; le mot lu « fera » (« feue ») est à la place où l'on attendrait « avec » ; le $x$ de « $xy$ » est récrit sur une autre lettre. Le feuillet s'arrête là, aux deux tiers de la hauteur ; le reste est blanc.}

\page{635}

Dans le cylindre le plan du cercle qui sert debase est celui
des $zy$. parconséquent \struck{l'arc} l'élément $ds$ de l'arc est l'élément
du même cercle. Pour avoir le cas de l'anneau il suffit donc
de faire $x$ constant dans l'équation du cylindre puisque cette
condition ne peut apartenir qu'aux seuls points dela circonférence
du cercle dont le plan est perpendiculaire a l'axe des $x$.
Lorsqu'il s'agit d'un cône le plan des $yz$ n'est plus lemême
que leplan du cercle qui lui sert de base. l'élément $ds$ del'arc
qui continue a être situé dans le plan de $yz$ ne se confond plus
avec l'élément dela base du cône il lui est au contraire
plus ou moins incliné \struck{jusqu'a ceque dans la plaque} en même tems
que son rayon de courbure augmente en sorte que dans la plaque
circulaire, qui est la limite du cône l'élément $ds$ devient l'élément
de l'ordonnée $dy$ en même tems qu'il est perpendiculaire a l'élément
du cercle extérieur dela plaque, cercle qui est la base de ce cône
infiniment \uncertain{surbaissé}.
\note{En haut à droite « 8 », sans soulignement, de sa main semble-t-il : sa pagination. Aucun numéro de la série à l'encre sur cette page : elle est le dos du feuillet qui porte « 311 » sur l'autre face, restée blanche (vue 634). Le haut de la vue montre, au-dessus de ce feuillet, le dos du feuillet de la vue 632, où son texte se lit à l'envers par transparence. Le « ne » de « ne peut » est repassé. La première lettre de « extérieur » est récrite.}

En prenant pour rayon la distance del'origine a un point
donné dela circonférence du cercle qui mesure la base du solide
dont on considère la surface et nommant ce rayon \uncertain{$r$}
on trouvera que si le rayon $z\,\text{cos}\,A$ \quad $A$ étant l'an
\note{Après « ce rayon », une ou deux petites lettres soulignées, lues avec doute $r$. La page s'arrête au milieu de la hauteur, sur « l'an », coupé en fin de ligne ; la suite n'est pas sur cette page.}

\page{636}

dont les rayons sont $r$ et $r + dr$ elle est comprise entre
les deux plans qui coupent les sphères dont les rayons sont
$r + dr$ suivant les cercles qui servent debases aux mêmes
calottes sphériques.
\note{En haut à droite, à l'encre, « 312 » ; à sa droite, un « 7 » traversé d'un trait oblique, le même signe qu'à la vue 632. La page commence au milieu d'une phrase (« dont les rayons sont ») : elle continue le bas de la vue 639, paginée 6, qui s'arrête sur « au lieu d'être comprise entre les deux calotte sphériques ». Elle ne suit pas la vue 635, qui s'arrête sur « l'an ».}

Les trois côtés de cette molécule sont $dr$, $r\,\text{sin}\,\theta\,d\varphi$, $r\,\text{cos}\,\theta\,d\theta$
les trois faces sont $r^2\,\text{sin}\,\theta\,\text{cos}\,\theta\,d\varphi\,d\theta$, $r\,\text{sin}\,\theta\,d\varphi\,dr$, et $r\,\text{cos}\,\theta\,d\theta\,dr$.
Pendant l'instant $dt$ la molécule reçoit par ces trois faces des quantités
de chaleur qui sont exprimées par les formules suivantes:
\[ -K\,r^2\,\text{sin}\,\theta\,d\varphi\,d\theta\,\text{cos}\,\theta.\frac{dv}{dr}.dt \]
\[ -K\,r\,\text{sin}\,\theta.d\varphi\,dr.\frac{dv}{d.r\,\text{sin}\,\theta}.\text{cos}\,\theta.dt \]
\[ -K\,r\,\text{cos}\,\theta\,d\theta\,dr.\frac{dv}{r\,\text{sin}\,\theta\,d\varphi}\,dt \]
les quantités de chaleur qui passent par les faces opposées sont
\[ \left\{-K\,r^2\,\text{sin}\,\theta\,\text{cos}\,\theta\,d\varphi\,d\theta\,\frac{dv}{dr} - K\,\text{sin}\,\theta\,\text{cos}\,\theta\,d\varphi\,d\theta\,d\Big(r^2\frac{dv}{dr}\Big)\right\} \]
\[ \left\{-K\,r\,\text{sin}\,\theta\,d\varphi\,dr\,\frac{dv}{dr\,\text{sin}\,\theta}\Big)\,\struck{\text{cos}\,\theta} - K\,r\,d\varphi\,d\theta.d.\Big(\text{sin}\,\theta\,\text{cos}\,\theta\,\frac{dv}{d(r\,\text{sin}\,\theta)}\Big)\right\} \]
\[ \left\{-K\,r\,\text{cos}\,\theta\,d\theta\,dr\,\frac{dv}{r\,\text{sin}\,\theta\,d\varphi} - K\,r\,d\theta\,dr\,\add{\text{cos}\,\theta}\,d.\Big(\frac{dv}{r\,\text{sin}\,\theta\,d\varphi}\Big)\right\} \]
\note{Les trois formules de gauche et les trois accolades de droite sont écrites côte à côte ; les mots « les quantités de chaleur qui passent par les faces opposées sont » sont disposés en colonne devant les accolades. Au bout de la première accolade, contre le bord du feuillet, une lettre coupée (« d… »). Dans la deuxième, une parenthèse fermante sans ouvrante après le premier dénominateur, et le « $\text{cos}\,\theta$ » qui suit est barré ; le $d\theta$ de « $d\varphi\,d\theta$ » est empâté. Le « $\text{cos}\,\theta$ » de la troisième est écrit au-dessus de la ligne.}

La première differentielle étant prise par rapport a \emph{r} seulement la
seconde demême par rapport a \ill{} et la troisième par rapport à $\varphi$.
En retranchant les quantités de chaleur qui entrent par chacune
des faces \struck{dela} de celles qui s'écoulent par les faces respectivement
opposées ontrouve l'équation
\note{Après « par rapport a », un signe court repassé à l'encre, que le sens voudrait être $\text{sin}\,\theta$ ; on ne le lit pas.}

\begin{align*}
r^2\,\text{sin}\,\theta\,\text{cos}\,\theta\,d\theta\,d\varphi\,dr\,\frac{dv}{dt}\,dt = \frac{K}{C.D} & \Big\{\text{sin}\,\theta\,\text{cos}\,\theta.d\varphi\,d\theta\,d\Big(r^2\frac{dv}{dr}\Big) + r\,d\varphi\,d\theta.d\Big(\text{sin}\,\theta\,\text{cos}\,\theta.\frac{dv}{d(r\,\text{sin}\,\theta)}\Big) \\
& + r\,\text{cos}\,\theta\,d\theta\,dr\,d\Big(\frac{dv}{r\,\text{sin}\,\theta\,d\varphi}\Big)\Big\}
\end{align*}
\[ \frac{dv}{dt} = \frac{K}{C.D}.\left\{\frac{1}{r^2}\,\frac{d\big(r^2.\frac{dv}{dr}\big)}{dr} + \frac{d\big(\text{sin}\,\theta\,\text{cos}\,\theta\,\frac{dv}{dr\,\text{sin}\,\theta}\big)}{r\,\text{sin}\,\theta\,\text{cos}\,\theta\,d\theta} + \frac{1}{r^2\,\text{sin}^2\theta}.\frac{d^2v}{d\varphi^2}\right\} \]
\[ = \frac{K}{C.D}\left\{\frac{d^2v}{dr^2} + \frac{2}{r}\,\frac{dv}{dr} + \frac{\text{cos}^2\theta}{r^2}.\frac{d^2v}{(d.\text{sin}\,\theta)^2} + \Big(\frac{\text{cos}\,\theta}{r^2\,\text{sin}\,\theta} - \frac{\text{sin}\,\theta}{r^2\,\text{cos}\,\theta}\Big)\frac{dv}{d\,\text{sin}\,\theta} + \frac{1}{r^2\,\text{sin}^2\theta}\,\frac{d^2v}{d\varphi^2}\right\} \]
\note{Première équation : le premier « $\text{sin}\,\theta$ » de l'accolade est précédé d'un trait noirci, et le groupe « $d(r^2$ » est couvert d'une tache d'encre ; le $dt$ qui suit $\frac{dv}{dt}$ est empâté. La lettre lue $D$ dans « $C.D$ » est un dôme sur une barre, comme celle que le lot 28 a relevée dans le brouillon sur la chaleur (vue 549) et qui pourrait être un $\Omega$. Les trois lignes s'étendent jusqu'au bord droit du feuillet.}

cette équation comparée a celle que j'ai trouvé d'une
autre manière présente cette difference que les termes
qui contiennent $\frac{d^2v}{(d\,\text{sin}\,\theta)^2}$ et $\frac{dv}{d\,\text{sin}\,\theta}$ \add{sont multipliés} \struck{doivent que} dans cette dernière
\add{par $\text{cos}\,\theta$} pour en trouver la raison il faut observer que l'on a multiplié
le terme $K\,r\,\text{sin}\,\theta\,d\varphi\,dr\,\frac{dv}{d(r\,\text{sin}\,\theta)}$ qui exprime la quantité de chaleur
qui traverse la face $r\,\text{sin}\,\theta\,d\varphi\,dr$ par $\text{cos}\,\theta$ parceque la direction $\frac{dv}{d\,\text{sin}\,\theta}$
fait l'angle $\theta$ avec la perpendiculaire au plan de cette face.
cela posé la chaleur qui traverse la molécule pour s'écouler $\times$
\marginal{$\times$ par la face opposée la traverse dans une direction oblique ainsi la difference $d\big(\text{sin}\,\theta\,\text{cos}\,\theta\,\frac{dv}{d\,\text{sin}\,\theta}\big)$ $K\,r\,d\varphi\,dr$ dela quantité de chaleur qui entre \uncertain{obliquement} \uncertain{et sort} par la face opposée ne mesure}

\emph{Retournez}
\note{La note marginale, écrite dans la marge de gauche en face des trois dernières lignes de l'équation et de tout ce paragraphe, continue la phrase du texte : le « $\times$ » qui clôt « pour s'écouler » l'appelle, et elle commence par le même signe. Sa formule est soulignée ; « par $\text{cos}\,\theta$ », au-dessus d'elle, appartient à la ligne du texte (« sont multipliés … par $\text{cos}\,\theta$ »). « obliquement » est écrit en abrégé et récrit ; « et sort » est repassé. Au pied de la marge, souligné, « Retournez » : la suite est au dos du feuillet (vue 637).}

\page{637}

plus la quantité de chaleur accumulée dans la molécule
mais celle qui seroit accumulée dans une autre molecule
que la chaleur traverseroit dans une direction perpendiculaire
au plan des deux faces opposées. Il est clair que
solidité dela première molecule est a la solidité
de la seconde $::\ \text{cos}\,\theta : 1$ parconséquent il
faut multiplier l'expression dela quantité de
chaleur qui seroit accumulée dans la molecule
traversée dans la direction perpendiculaire
quantité qui seroit $d\big(\text{sin}\,\theta\,\text{cos}\,\theta\,\frac{dv}{d(r\,\text{sin}\,\theta)}\big)$ par
$\text{cos}\,\theta$ et on aura ainsi l'expression dela
quantité dechaleur qui est accumulée dans
la molecule qui est traversée obliquement
par le flux de chaleur. alors les équations
seront identiques.
\note{Dos du feuillet de la vue 636, qui porte au pied de sa marge « Retournez » : la première ligne continue la note marginale de ce recto (« ne mesure plus la quantité de chaleur accumulée dans la molécule »). Aucun numéro de la série à l'encre sur cette face. Le « 8 » qui paraît en haut à droite de la vue est sur une étroite bande de papier visible au-dessus de ce feuillet, qui semble être le haut de la page de la vue 635 ; il n'appartient pas à cette page. « qui est traversée » : « est » est récrit. Le texte s'arrête au milieu du feuillet sur « seront identiques. », suivi d'une tache ; le reste est blanc, et le recto s'y lit par transparence.}

\page{639}

Les trois côtés de cet élément sont $dr$, $r\,\text{sin}\,\theta\,d\varphi$ et $r\,d\theta$
les trois faces sont $r^2\,\text{sin}\,\theta\,d\varphi\,d\theta$, $r\,\text{sin}\,\theta\,d\varphi\,dr$ et $r\,d\theta\,dr$,
Pendant l'instant $dt$ la molécule reçoit par ces trois
faces des quantités de chaleur qui seront \struck{représentent} exprimées
par les formules.
\note{En haut à droite, à l'encre, « 312 » barré d'un trait horizontal et, à sa droite, « 6 » barré d'un trait oblique. L'autre face de ce feuillet (vue 638) est blanche et porte « 313 » : le feuillet a d'abord été numéroté 312 sur cette face, puis 313 sur l'autre ; c'est la vue 636 qui porte 312. Au-dessus de ce feuillet, sur une étroite bande de papier, le même « 8 » qu'à la vue 637. « cet » est récrit sur « cette ».}

\[ -K\,r^2\,\text{sin}\,\theta\,d\varphi\,d\theta.\frac{dv}{dr}.dt \]
\[ -K\,r\,\text{sin}\,\theta\,d\varphi\,dr.\frac{dv}{r\,d\theta}\,dt \]
\[ -K\,r\,d\theta\,dr\,\frac{dv}{r\,\text{sin}\,\theta\,d\varphi}.dt \]
les quantités de chaleurs qui passent par les faces opposées sont
\[ \left\{-K\,r^2\,\text{sin}\,\theta\,d\varphi\,d\theta\,\frac{dv}{dr} - K\,\text{sin}\,\theta\,d\varphi\,d\theta\,d\Big(r^2\frac{dv}{dr}\Big)\right\}dt \]
\[ \left\{-K\,r\,\text{sin}\,\theta\,d\varphi\,dr\,\frac{dv}{r\,d\theta} - K\,r\,d\varphi\,dr.d\Big(\text{sin}\,\theta\,\frac{dv}{r\,d\theta}\Big)\right\}dt \]
\[ \left\{-K\,r\,d\theta\,dr\,\frac{dv}{r\,\text{sin}\,\theta} - K\,r\,d\theta\,dr\,d.\frac{dv}{r\,\text{sin}\,\theta\,d\varphi}\right\} \]
\note{Même disposition qu'à la vue 636 : trois formules à gauche, trois accolades à droite, les mots « les quantités de chaleurs qui passent par les faces opposées sont » en colonne entre les deux. Dans la première formule de gauche, le $d$ du numérateur est empâté ; les $\theta$ des dénominateurs de la deuxième et de la troisième sont repassés. Le premier dénominateur de la troisième accolade s'arrête sur « $r\,\text{sin}\,\theta$ », sans $d\varphi$ : ainsi sur le feuillet.}

la première differentielle étant prise par rapport a $r$
\add{seulement} la seconde \struck{differe} \add{demême} par rapport a $\theta$ \struck{et} et la troisième
par rapport a $\varphi$.
\note{« seulement » est écrit dans la marge de gauche, surmonté d'un mot court (« demême » ?) ; « demême » est écrit au-dessus de « differe » barré.}

En retranchant les quantités de chaleur qui entrent par
chacune des faces de \struck{l'élément} la molécule \struck{sphériques} \add{dela sphère} de
celles qui s'écoulent par les faces respectivement opposées
on trouve l'équation
\[ r^2\,\text{sin}\,\theta\,d\theta\,d\varphi\,dr\,\frac{dv}{dt}\,dt = \frac{K}{C.D}.\left\{\text{sin}\,\theta\,d\varphi\,d\theta\,d\Big(r^2\frac{dv}{dr}\Big) + r\,d\varphi\,dr\,d\Big(\text{sin}\,\theta\,\frac{dv}{r\,d\theta}\Big) + \frac{r\,d\theta\,dr}{r\,\text{sin}\,\theta}\,d.\frac{dv}{d\varphi}\right\}dt \]
\[ \frac{dv}{dr} = \frac{K}{C\,D}.\left\{\frac{\ill{}}{r^2}\,\frac{d.\big(r^2\frac{dv}{dr}\big)}{dr} + \frac{d\big(\text{sin}\,\theta\,\frac{dv}{d\theta}\big)}{r^2\,\text{sin}\,\theta\,d\theta} + \frac{1}{r^2\,\text{sin}^2\theta}.\frac{d^2v}{d\varphi}\right\} \]
\[ = \frac{K}{C.D}\left\{\frac{d^2v}{dr^2} + \frac{2}{r}\,\frac{dv}{dr} + \frac{d^2v}{r^2\,d\theta^2} + \frac{\text{cos}\,\theta}{r\,\text{sin}\,\theta}.\frac{dv}{r\,d\theta} + \frac{1}{r^2\,\text{sin}^2\theta}.\frac{d^2v}{d\varphi^2}\right\} \]
\note{Le « $dt$ » qui ferme l'accolade de la première équation est écrit sous la ligne, au bord droit. « $\frac{dv}{dr}$ » en tête de la deuxième ligne : ainsi sur le feuillet, où le calcul attend $\frac{dv}{dt}$. Le numérateur de la première fraction de l'accolade est noirci d'une tache ; le sens attend 1. Le dernier terme de cette ligne porte un petit 2 au-dessus du $d$ du numérateur et aucun exposant au dénominateur. Le $d\theta^2$ de la troisième ligne est repassé. Pour « $C.D$ », voir la note de la vue 636.}

Si au lieu de considérer le solide conique qui s'appuie sur
\struck{la} calotte sphérique on veut s'occuper du \add{simple} cône dont la base
est \struck{la la plaque circulaire} est l'intersection dela sphère et
du plan qui passe par le parallèle donné on aura a
examiner au lieu de la molecule \struck{sphérique} \add{à la sphère} la molecule \struck{conique}
\struck{qui} du cône et celle-ci diffère \struck{en} dela première en \uncertain{ce}
point qu'au lieu d'être comprise entre les deux calotte sphériques
\note{La page s'arrête au bas du feuillet sur « calotte sphériques ». La phrase continue en tête de la vue 636 (« dont les rayons sont $r$ et $r + dr$ elle est comprise entre les deux plans … »), qui porte la pagination 7 quand celle-ci porte 6 : les deux feuillets sont reliés dans l'ordre inverse de sa pagination. Cette page traite la molécule de la sphère (côtés $dr$, $r\,\text{sin}\,\theta\,d\varphi$, $r\,d\theta$) ; la vue 636 reprend le même calcul pour la molécule du cône, où le côté est $r\,\text{cos}\,\theta\,d\theta$ et la variable $\text{sin}\,\theta$.}

\page{640}

\[ \Big(\frac{dz}{dx}\Big)^2 = \frac{\text{cos}^2\theta\,\text{cos}^2\varphi}{r^2}\,\Big(\frac{dz}{d\theta}\Big)^2 - \frac{2\,\text{cos}\,\theta\,\text{cos}\,\varphi\,\text{sin}\,\varphi}{r^2\,\text{sin}\,\theta}.\frac{dz}{d\theta}.\frac{dz}{d\varphi} + \frac{\text{sin}^2\varphi}{r^2\,\text{sin}^2\theta}.\Big(\frac{dz}{d\varphi}\Big)^2 \]
\begin{align*}
\frac{d^2z}{dy^2} = {} & \frac{d^2z}{d\theta^2}\,\frac{\text{cos}^2\theta\,\text{sin}^2\varphi}{r^2} + \frac{d^2v}{d\theta\,d\varphi}.\frac{2\,\text{cos}\,\theta\,\text{sin}\,\varphi\,\text{cos}\,\varphi}{r^2\,\text{sin}\,\theta} + \frac{d^2v}{d\varphi^2}\,\frac{\text{cos}^2\varphi}{r^2\,\text{sin}^2\theta} \\
& - \frac{dz}{d\theta}\left\{\frac{2\,\text{sin}\,\theta\,\text{cos}\,\theta\,\text{sin}^2\varphi}{r^2} - \frac{\text{cos}^2\varphi\,\text{cos}\,\theta}{r^2\,\text{sin}^2\theta}\right\}. \\
& - \frac{dz}{d\varphi}\left\{\frac{\ill{}\,\text{sin}\,\varphi\,\text{cos}\,\varphi}{r^2\,\text{sin}^2\theta} + \frac{\text{cos}^2\theta\,\text{sin}\,\varphi\,\text{cos}\,\varphi}{r^2\,\text{sin}^2\theta} + \frac{\text{sin}\,\varphi\,\text{cos}\,\varphi}{r^2}\right\}
\end{align*}
\note{En haut à droite, à l'encre, « 314 » ; aucune pagination de sa main. La page commence sans phrase d'introduction, sur le carré de $\frac{dz}{dx}$ ; le dernier facteur de la première ligne touche le bord droit du feuillet. Elle écrit $v$ dans deux numérateurs de $\frac{d^2z}{dy^2}$ où les autres portent $z$ : ainsi sur le feuillet. Le premier facteur du premier numérateur de la dernière accolade est couvert d'une tache.}

\emph{Recherches sur les \struck{surfa} cones}

S'il s'agit d'un cône \add{l'élément de} la surface parallèle a la base
est une portion de secteurs comprise entre deux cercles concentriques.
les rayons de ces cercles sont $r\,\text{sin}\,\theta$ et \struck{$r\,\text{cos}.d\theta\,\text{sin}(\theta - d\theta)$} \add{$r'\,\text{sin}(\theta - d\theta)$}
$r$ est le \struck{rayon} côté du cône qui a pour base le cercle
extérieur $\theta$ est l'angle que ce côté fait avec l'axe
du cône parconséquent on a l'axe
\[ z = r\,\text{cos}\,\theta = r'\,\text{cos}(\theta - d\theta) \qquad r' = \frac{r\,\text{cos}\,\theta}{\text{cos}(\theta - d\theta)} \]
\[ r'\,\text{sin}(\theta - d\theta) = r\,\text{cos}\,\theta\,\frac{(\text{sin}\,\theta\,\text{cos}\,d\theta - \text{cos}\,\theta\,\text{sin}\,d\theta)}{\text{cos}\,\theta\,\struck{\text{cos}\,d}\theta + \text{sin}\,\theta\,\text{sin}\,d\theta} \]
en négligeant les termes du second ordre $= \frac{r\,\text{cos}\,\theta\,(\text{sin}\,\theta - \text{cos}\,\theta\,d\theta}{\text{cos}\,\theta + \text{sin}\,\theta\,d\theta}$
\note{Le titre, en tête de la seconde partie de la page, est souligné d'un long trait. Devant « $r' =$ », une tache d'encre. Au-dessus du $r$ de « $r\,\text{cos}\,\theta$ », après le premier signe $=$ de la deuxième formule, un petit signe récrit ; on lit $r$. Au dénominateur, le « cos d » est noirci ; la parenthèse ouverte après $\text{cos}\,\theta$ dans la dernière fraction n'est pas fermée.}

\[ r\,\text{sin}\,\theta - \frac{r\,\text{cos}\,\theta\,(\text{sin}\,\theta - \text{cos}\,\theta.d\theta)}{\text{cos}\,\theta + \text{sin}\,\theta\,d\theta} = \frac{r\,\text{sin}\,\theta\,(\text{cos}\,\theta + \text{sin}\,\theta\,d\theta) - r\,\text{cos}\,\theta\,(\text{sin}\,\theta - \text{cos}\,\theta\,d\theta)}{\text{cos}\,\theta + \text{sin}\,\theta\,d\theta} \]
en effaçant cequi se détruit
et négligeant $d\theta\,\text{sin}\,\theta$ vis avis de $\text{cos}\,\theta$ on a
$\frac{r^2}{2}\,\frac{d\varphi\,d\theta\,\text{sin}\,\theta}{\text{cos}\,\theta}$ pour la surface cherchée que traverse la chaleur
dans la direction $\frac{dv}{dr}$. par l'élément compris entre $\frac{r\,d\theta}{\text{cos}\,\theta}$ et $dr$
dans la direction $\frac{dv}{d\varphi}$. $dr\,\frac{r\,d\theta}{\text{cos}\,\theta}$ \quad $\frac{d\uncertain{v}}{d\theta}\,r\,\text{sin}\,\theta\,d\varphi\,dr$
\note{Dans « $\frac{r^2}{2}\,\frac{d\varphi\,d\theta\,\text{sin}\,\theta}{\text{cos}\,\theta}$ », le « $\text{sin}\,\theta$ » est écrit en petit au-dessus de la ligne, après $d\theta$. À la dernière ligne, le numérateur lu $dv$ est empâté et le $\theta$ de « $\text{sin}\,\theta$ » récrit. La page s'arrête là, aux trois quarts du feuillet ; le reste est blanc.}

\end{document}
